% Source: https://artofproblemsolving.com/community/c3h1386939_mathcounts
% Stable identity: aops-topic-1386939-q02
Work from the bottom: if $n\equiv0\pmod3$ there is no $3^0$ term; if $n\equiv1$, use $+1$; if $n\equiv2$, use $-1$. Then divide $n\mp1$ by $3$ and repeat (this is \emph{balanced ternary}, and the representation is unique). For example $71\equiv2$, so $71=-1+72=-1+3\cdot24$; then $24=3\cdot8$, $8\equiv2$ so $8=-1+9$. Altogether $71=-1+3(3(-1+9))=-1-9+81$. The results:
\begin{align*}
45&=81-27-9, & 52&=81-27-3+1,\\
63&=81-27+9, & 71&=81-9-1,\\
89&=81+9-1, & 115&=81+27+9-3+1,\\
296&=243+81-27-1. &&
\end{align*}
